\section{结论}

本文没有把“心学突围、朝廷收编”当成一个已测得的历史事实，而是把它拆成关系构成与时间序列：思想相关组的显式关系中，师生同门占30.0\%，权力关系占58.9\%，显示传播与权力接合同时发生。

时间线上，1508--1528年集中了形成与传授事件，1579与1593年的书院节点进入政治斗争类别；这支持“叙事从个人悟道转向制度冲突”的语料内判断。

可检验的下一步，是把原典、奏疏和书院档案加入独立语料，重标关系类别并比较同一批人物。若组间构成和时间转折仍然出现，解释框架才有资格越过这一本通俗史著。
